\section{In-Context Learning}
\label{sec:in_context_learning}

In-context learning is a paradigm that allows language models to learn a task based on a few example inputs and outputs also called demonstrations~\cite{brown_language_2020,dong_survey_2023}. 
In the context of this paper, we combine the closed prompt `Generated Knowledge` and the open prompt with demonstrations to analyze the impact of demonstrations on the performance of ChatGPT for Product Information Extraction. First, %ly, 
we explain how the demonstrations are selected and combined with the prompts. %Secondly, 
Then, we present the results of our experiments.

\vspace{.2cm}\noindent\textbf{Demonstration Selection.}
The source for the used demonstrations is the reduced training set. Demonstrations are useful if they are semantically similar to the final task~\cite{liu_what_2022}. Since we know the product category of the input product title and assume that product offers with the same product category are semantically similar, the selected demonstrations are drawn from the same product category as the product offer for which we run the product information extraction. Both the closed and the open prompt are combined with one, three and six demonstrations. For the prompts with three and six demonstrations, two-thirds of the demonstrations are product offers where the value of the attribute/value pair is contained in the title and one-third are product offers where the value of the attribute/value pair is not contained in the title. This selection of demonstrations covers different possible structures of the task~\cite{levy_diverse_2022}. All attribute values of the demonstrations are manually checked and normalized like the examples listed in Table \ref{tab:value_normalization}.
%similar to the ground truth attribute values such as \emph{"8 hour"} and \emph{"13.3-inch"} which are normalized to \emph{"8 hours"} and \emph{"13.3 inches"}, respectively.
%Like in the ground truth attribute values like \emph{``8 hour``} and \emph{``13.3-inch``} are normalized to \emph{``8 hours``} and \emph{``13.3 inches``}.

\vspace{.2cm}\noindent\textbf{Combination of Prompts and Demonstrations.} For the in-context learning prompts, we first provide the task description and add the task demonstrations. Each task demonstration consists of an example input and an example output. Afterwards, we repeat the task description and add the task input.
The prompt examples (e) and (f) in Figure \ref{fig:few_shot_cropped} illustrate closed and open prompts with demonstrations.

%\begin{figure}[ht]
%\includegraphics[width=.45\textwidth]{figures/gk_few_shots_cropped_V2.pdf}
%\caption{This example prompt utilizes demonstrations to guide the closed product information extraction.}
%\label{fig:gk_few_shots_cropped}
%\end{figure}


%\begin{figure}[ht]
%\includegraphics[width=.45\textwidth]{figures/open_few_shots_cropped_V2.pdf}
%\caption{This example prompt utilizes demonstrations to guide the open product information extraction.}
%\label{fig:open_few_shots_cropped}
%\end{figure}

\begin{figure}[ht]
\includegraphics[width=.475\textwidth]{figures/few_shot_cropped.pdf}
\caption{This example prompt utilizes demonstrations to guide the open product information extraction.}
\label{fig:few_shot_cropped}
\end{figure}

\begin{table}[]
\setlength{\tabcolsep}{4pt}
\caption{Precision (P), Recall (R), F1, $\Delta$ F1 and Cost (\textcent) per Title of the prompts with one, three and six demonstrations (Shots) and fully trained baselines}
\label{tab:few_shot_results}
\scalebox{0.925}{\begin{tabular}{@{}ll|ccccc@{}}
\toprule
Config. & \# Shots & P & R & F1 & $\Delta$F1 over & Cost (\textcent)                   \\ 
 &  &  &  &  & zero-shot  &  per Title                    \\ \midrule
Closed        & Zero  & 87.79                 & 87.99                 & 87.89                  & \multicolumn{1}{c}{-}        & \textbf{0.0242}                     \\
Closed        & One   & 90.99                 & 93.30                  & 92.13                  & +4.24                         & 0.0482                     \\
Closed        & Three & 91.17                 & 95.38                 & 93.23                  & +5.34                         & 0.0705                     \\
Closed        & Six   & 93.17                 & \textbf{95.94}                 & \textbf{94.51}                  & +6.62                         & 0.1037                     \\ \midrule
Open          & Zero  & 86.92                 & 43.46                 & 57.95                  & -29.94                       & 0.0305                     \\
Open          & One   & \textbf{95.65}                 & 86.37                 & 90.77                  & +2.88                         & 0.0679                     \\
Open          & Three & 91.08                 & 81.03                 & 83.92                  & -3.97                        & 0.0910                     \\
Open          & Six   & 91.48                 & 82.77                 & 85.73                  & -2.16                        & 0.1203   \\ \midrule
\multicolumn{2}{l|}{NER (fully trained)}      & 90.80                  & 84.3                 & 87.43                  &        \multicolumn{1}{c}{-}  & \multicolumn{1}{c}{-}   \\
\multicolumn{2}{l|}{AVEQA (fully trained)} & \textbf{99.52} & \textbf{99.28} & \textbf{99.41} &\multicolumn{1}{c}{-} &\multicolumn{1}{c}{-} 
\\ \bottomrule
\end{tabular}}
\end{table}

\vspace{.2cm}\noindent\textbf{Results.}
The results with in-context learning in Table \ref{tab:few_shot_results} show that demonstrations improve the extraction of attribute/value pairs of closed and open prompts. A single example (one-shot) is already sufficient to increase the closed extraction results by 4.24\% in terms of F1. Additional demonstrations add another 1-2\% in F1, with a maximal F1 of 94.51 achieved with (only) six samples.

The open extraction benefits the most from the demonstrations. A single example is sufficient to increase the performance of ChatGPT by almost 33\% in terms of F1. We noticed that the model picks up on the structure of the demonstrations and uses the demonstration attribute/value pairs to extract the attribute/value pairs from the task input. Looking at the example prompt (f) in Figure \ref{fig:few_shot_cropped}, for example, the model learns from the task demonstration that the attribute value \emph{"10 MP"} has to be extracted as \emph{"Resolution: 10 MP"}, rather than \emph{"Megapixels: 10 MP"} as it can happen if no demonstration is provided.
%ChatGPT's F1 increases if an open prompt with one demonstration is used but this comes at the cost of a high loss in recall compared to closed extraction.
In our experiments, open extraction prompts with three and six demonstrations resulted in a worse performance than experiments with only one demonstration. We assume that this behaviour is caused by the longer prompts but additional experiments are necessary to verify this hypothesis.

\vspace{.2cm}\noindent\textbf{Cost of the extraction.}
%\rs{What do you think about explicitly expressing here the trade-off between high-cost and high-quality and the way they are expressed in open/closed extraction? I think it can make the final paragraph more general; yet, still expressing what you wanted in terms of the cost.} 
On the one hand, the closed extraction few-shot prompts are 25\% cheaper per API call and exhibit an 8\% higher F1 recall than the open extraction few-shot prompts. On the other hand, open extraction prompts can extract multiple attributes, which makes them cheaper than closed extraction prompts if two or more attributes have to be extracted. 

%Todo Alex (Done): Shorten list of authors in the references using "et. al."